\section{Conclusion}\label{sec:conclusion}

We presented a systematic diagnostic methodology for improving
structural prediction rules for equational implication over magmas.
By treating each rule as a binary classifier and measuring its
precision across multiple test splits, we identified and corrected
rules that were anti-correlated with the ground truth on adversarial
instances.  Our main theoretical contributions---the projection
collapse lemmas (Theorems~\ref{thm:left-proj}
and~\ref{thm:right-proj})---are elementary but practically valuable,
handling a class of implications that purely structural heuristics
miss.

The refined predictor achieves $79.1\%$ overall accuracy (up from
$69.8\%$) and $63.5\%$ on the hardest split (up from $52.75\%$),
while fitting within a $7.8$\,KB text cheatsheet.  The remaining
errors ($146$ out of $400$ on \textsc{Hard3}) are dominated by
cases requiring algebraic reasoning beyond syntactic features.

Several directions for future work suggest themselves.

\begin{enumerate}[label=(\roman*)]
\item \textbf{Extended collapse classification.}  Theorems~\ref{thm:left-proj}
  and~\ref{thm:right-proj} handle the simplest projection-forcing
  equations.  A systematic classification of all equations in~$\Eq$
  that force a projection algebra (or more generally, that force the
  magma into a small number of isomorphism classes) would yield
  additional implication rules.

\item \textbf{Algebraic rewriting integration.}  Simple rewriting
  chains---such as detecting idempotent laws ($x * x = x$) as
  consequences---could be encoded as additional prediction rules.
  The challenge is to identify rewriting patterns that are both
  common enough to improve accuracy and simple enough to fit in a
  compact cheatsheet.

\item \textbf{Counterexample tables.}  The $16$ magmas of order~$2$
  (or a curated subset of the $524$ critical magmas
  from~\cite{bolan2025equational}) could be included in the cheatsheet
  as a verification step.  A counterexample refuting $E_1 \models E_2$
  provides a definitive \textsc{false} prediction and could eliminate
  many of the residual false positives from contradiction motifs.

\item \textbf{Stone pairing features.}  The geometric structure
  revealed by Berlioz and Melli\`es~\cite{berlioz2026latent} suggests
  that approximate Stone pairing values, computed on a few canonical
  magmas, could serve as continuous features for prediction.  This
  would bridge the gap between our discrete structural rules and the
  continuous geometry of the implication lattice.

\item \textbf{Confidence calibration.}  Our current predictor outputs
  binary \textsc{true}/\textsc{false} predictions.  Assigning
  calibrated confidence scores---based on the number and reliability
  of rules that fire---could improve log-loss performance, which is
  the primary competition metric.
\end{enumerate}
